\section{Distribución $t$}
\label{sec:distribucion-t}

La distribuci\'on $t$ aparece cuando estandarizamos una media muestral usando la
desviaci\'on est\'andar muestral (que es aleatoria) en lugar de la poblacional.

\begin{definicion}[Distribuci\'on $t$ de Student]
Sean $Z \sim N(0,1)$ y $\chi^2_\nu$ una variable chi-cuadrada con $\nu$ grados de
libertad, independientes. Entonces
\begin{align}
 \label{eq:3.2.12}
 T = \frac{Z}{\sqrt{\chi^2_\nu / \nu}}
\end{align}
sigue una distribuci\'on \emph{$t$ de Student} con $\nu$ grados de libertad. Escribimos
$T \sim t_\nu$.
\end{definicion}

{Propiedades}
\begin{itemize}
 \item La densidad de $T$ es
 \begin{align}
  \label{eq:3.2.13}
  f(t) = \frac{\Gamma\!\left(\frac{\nu+1}{2}\right)}{\sqrt{\nu\pi}\,\Gamma\!\left(\frac{\nu}{2}\right)}
  \left(1 + \frac{t^2}{\nu}\right)^{-(\nu+1)/2}, \quad t \in \mathbb{R}.
 \end{align}
 \item $E(T) = 0$ para $\nu > 1$.
 \item $\text{Var}(T) = \nu/(\nu - 2)$ para $\nu > 2$.
 \item Cuando $\nu \to \infty$, $T \xrightarrow{d} N(0,1)$. Para $\nu$ peque\~no,
 $T$ tiene colas m\'as pesadas que la normal.
\end{itemize}

\begin{observacion}
La distribuci\'on $t$ surge al estimar la media $\mu$ con $\sigma$ desconocida:
si $X_1, \ldots, X_n \sim N(\mu, \sigma^2)$ iid y $S^2$ es la varianza muestral,
entonces
\begin{align}
 \label{eq:3.2.14}
 \frac{\bar{X} - \mu}{S/\sqrt{n}} \sim t_{n-1}.
\end{align}
Este es el \emph{estad\'istico t} que da nombre a la prueba-$t$ (secci\'on \ref{sec:prueba-t-media}). Esta
es precisamente la descomposici\'on $T = Z/\sqrt{\chi^2_{n-1}/(n-1)}$ obtenida a
partir del Teorema de Fisher en la secci\'on anterior.
\end{observacion}

\subsubsection{Intervalos de confianza con $\sigma$ desconocida: el caso de muestras
peque\~nas}

Cuando $\sigma$ es desconocida (el caso pr\'actico m\'as com\'un), el intervalo de
confianza para $\mu$ debe usar los cuantiles de la distribuci\'on $t$ en lugar de los
de la Normal est\'andar.

\begin{teorema}[Intervalo de confianza para $\mu$ con $\sigma$ desconocida]
 \label{eq:5.4.1}
Sea $X_1, \ldots, X_n \sim N(\mu, \sigma^2)$ iid con $\sigma^2$ desconocida. Un
intervalo de confianza del $(1-\alpha)\times 100\%$ para $\mu$ es
\begin{align}
 \bar{X} \pm t_{n-1,\,\alpha/2} \cdot \frac{S}{\sqrt{n}},
\end{align}
donde $t_{n-1,\,\alpha/2}$ es el cuantil superior $\alpha/2$ de $t_{n-1}$.
\end{teorema}

\begin{observacion}[Por qu\'e importa m\'as en muestras peque\~nas]
Como $t_{n-1,\,\alpha/2} > z_{\alpha/2}$ para todo $n$ finito (las colas de $t$ son
m\'as pesadas), el intervalo basado en $t$ es siempre \emph{m\'as ancho} que el
intervalo basado en $Z$ que se usar\'ia si $\sigma$ fuera conocida. Esta diferencia es
sustancial para $n$ peque\~no (e.g., $t_{9,0.025}\approx 2.262$ vs.
$z_{0.025}=1.96$, una diferencia de m\'as del $15\%$) y se vuelve despreciable
cuando $n \to \infty$, ya que $t_\nu \xrightarrow{d} N(0,1)$.
\end{observacion}

\begin{ejemplo}
 \label{exmp:5.4.1}
Una muestra de $n=9$ mediciones de resistencia de un material produce
$\bar{X}=48.5$ MPa y $S=3.2$ MPa. Construya un intervalo de confianza del $95\%$
para la resistencia media $\mu$, usando $t_{8,0.025}\approx 2.306$.
\end{ejemplo}

\begin{solucion}
El error est\'andar es $S/\sqrt{n} = 3.2/3 \approx 1.067$. El intervalo es
\begin{align}
 \bar{X} \pm t_{8,0.025}\cdot\frac{S}{\sqrt{n}} = 48.5 \pm 2.306\cdot 1.067
 \approx 48.5 \pm 2.46 = [46.04,\ 50.96] \text{ MPa}.
\end{align}
Si en cambio hubi\'eramos usado (incorrectamente) el cuantil Normal
$z_{0.025}=1.96$, el intervalo habr\'ia sido $48.5 \pm 2.09 = [46.41, 50.59]$: m\'as
angosto de lo que la incertidumbre real de una muestra de solo $9$ observaciones
justifica.
\end{solucion}


